\begin{definition}
\label{definition-regular-ideal}
Let $R$ be a ring and let $I \subset R$ be an ideal.
\begin{enumerate}
\item We say $I$ is a {\it regular ideal} if for every
$\mathfrak p \in V(I)$ there exists a $g \in R$, $g \not \in \mathfrak p$
and a regular sequence $f_1, \ldots, f_r \in R_g$ such that $I_g$
is generated by $f_1, \ldots, f_r$.
\item We say $I$ is a {\it Koszul-regular ideal} if for every
$\mathfrak p \in V(I)$ there exists a $g \in R$, $g \not \in \mathfrak p$
and a Koszul-regular sequence $f_1, \ldots, f_r \in R_g$ such that $I_g$
is generated by $f_1, \ldots, f_r$.
\item We say $I$ is a {\it $H_1$-regular ideal} if for every
$\mathfrak p \in V(I)$ there exists a $g \in R$, $g \not \in \mathfrak p$
and an $H_1$-regular sequence $f_1, \ldots, f_r \in R_g$ such that $I_g$
is generated by $f_1, \ldots, f_r$.
\item We say $I$ is a {\it quasi-regular ideal} if for every
$\mathfrak p \in V(I)$ there exists a $g \in R$, $g \not \in \mathfrak p$
and a quasi-regular sequence $f_1, \ldots, f_r \in R_g$ such that $I_g$
is generated by $f_1, \ldots, f_r$.
\end{enumerate}
\end{definition}